\section{Year-wise Description of Growth Rates in Bangladesh (1951-2023)}

\subsection{Plots}

\begin{figure}[H]
    \centering
    \includegraphics[width=0.8\linewidth]{TOTAL GR.png}
     
    \label{fig:placeholder}
\end{figure}

\begin{figure}[H]
    \centering
    \includegraphics[width=0.8\linewidth]{GR(WP VS DP).png}
     
    \label{fig:placeholder}
\end{figure}

\subsection{Interpretation of the plot for working age population growth rates vs dependents(child+adult) growth rates
}

\textcolor{purple}{1950s--1960s}
\begin{itemize}[leftmargin=*]
    \item \textbf{Child + Adult Dependent Growth Rate (blue):} Rising steadily, peaking around 1964--65 at nearly 1.7\%. This shows a rapid increase in population, especially among children, consistent with high fertility and declining infant mortality after independence from British India.
    \item \textbf{Working Population Growth Rate (orange):} Stable around 2.3--2.5\%, reflecting a young but expanding labor force.
\end{itemize}

\textcolor{purple}{1970--1975}
\begin{itemize}[leftmargin=*]
    \item \textbf{1971 Liberation War:} Both WP and DP growth rates plummeted sharply around 1971--72 due to large-scale deaths, displacement, and destruction during the Bangladesh Liberation War.
    \begin{itemize}
        \item DP GR (blue) dropped close to 0.9\%, while WP GR (orange) also dipped to around 1.6\%.
    \end{itemize}
    \item \textbf{1974 Famine:} Population growth suffered another setback in the mid-1970s with widespread famine. This explains the fluctuations and slower recovery in both dependent and working population growth during 1973--75.
\end{itemize}

\textcolor{purple}{Late 1970s--1980s}
\begin{itemize}[leftmargin=*]
    \item \textbf{Dependents (blue):} After recovering from the famine, the dependent population growth stabilized around 1.2--1.4\%, but started declining after 1985.
    \item \textbf{Working Population (orange):} Increased strongly to nearly 2.8--2.9\% by the mid-1980s, reflecting the demographic momentum.
    \item This period marks the early demographic transition---fertility decline begins, but the working population expands rapidly.
\end{itemize}

\textcolor{purple}{1990s}
\begin{itemize}[leftmargin=*]
    \item \textbf{Dependents:} Show continuous decline, falling below 0.5\% by 1995--2000, signaling fertility reduction and smaller family sizes.
    \item \textbf{Working Population:} Peaks near 3.0\% in the early 1990s, creating a demographic dividend opportunity.
    \item \textbf{1999--2000 (Sierra Leone War):} Bangladesh participated in UN peacekeeping operations in Sierra Leone during this period. The domestic demographic situation showed:
    \begin{itemize}
        \item Low dependent growth (near zero).
        \item Working population still strong ($\approx 2.5\%$).
    \end{itemize}
    This reflects Bangladesh’s growing role in the international community, supported by a robust labor force.
\end{itemize}

\textcolor{purple}{2000s}
\begin{itemize}[leftmargin=*]
    \item \textbf{Dependents:} Growth rate becomes nearly stagnant ($\sim 0.2$--0.3\%), showing population stabilization in child cohorts.
    \item \textbf{Working Population:} Declines gradually from $\sim 2.5\%$ (2000) to 1.5\% (2010), reflecting slowing labor force growth as fertility drops further.
\end{itemize}

\textcolor{purple}{2010s}
\begin{itemize}[leftmargin=*]
    \item \textbf{Dependents:} Near zero or negative growth around 2012--2016, indicating aging population trends and very low fertility.
    \item \textbf{Working Population:} Continues to decline, dipping below 1.0\% by 2017--18.
    \item \textbf{2019 Dengue Outbreak:} A major dengue outbreak in Bangladesh coincides with the period when both WP and DP growth rates were already low. While not directly visible in macro growth rates, it reflects the public health challenges faced during a slowing demographic momentum.
\end{itemize}

\textcolor{purple}{2020--2023}
\begin{itemize}[leftmargin=*]
    \item \textbf{COVID-19 Pandemic and Recovery:} Both series show a slight rebound after 2020.
    \item \textbf{Dependents:} Growth rate rises modestly ($\approx 0.3$--0.4\%), possibly due to short-term fertility and mortality shifts.
    \item \textbf{Working Population:} Recovers from its lowest point ($<1\%$ in 2019) to $\sim 1.6\%$ in 2023, suggesting a temporary demographic boost post-pandemic.
    \item \textbf{Migration Effects:} Bangladesh has a high rate of international labor migration (Middle East, Southeast Asia). During COVID, many migrants returned home, boosting the domestic working population temporarily. As migration resumed in 2021--22, the labor force remained large but more stable.
\end{itemize}
